% Pursuing Stacks (cotes 134-2 to 134-8): archive cote 134-3, site page 15 (PDF page 16), typescript page 15.
% typescript. Draft compared: ps2.tex:669-728.
% AI-transcribed by Claude Opus 5.5 (claude-opus-5-5, adaptive thinking, effort high), Message Batches API, batch ps1 (msgbatch_01Rv5hCaNTtUt4JHx1yqoHPL), 2026-11-12; not yet checked by eye.
\scanpage{134-3/15}
\mspage{15}

the category  (Ss sets) (and hence \over{\unsure{the}}{any} induced categor\over{ies}{y}) \over{are}{is} already
a topos,  namely equivalent to the category of sheaves on the same for
the canonical topology,  whereas in the category (Cat)  this is not so.

These reflections suggest that in most if not all categories of
models encountered for describing usual homotopy types,  there is a \margin{\unsure{5.3.} this \unsure{presumption} is false actually}
natural structure of a site on the model category M (\uline{presumably} the one
corresponding to the canonical topology(\over{*}{15})), with the property that the
homotopy type of any object X in M is \xout{also} canonically isomorphic to the
(pro)homotopy type of the induced site $M_{/X}$ ,  or what amount\hand{s} to the same,
of the corresponding topos
\[ (M_{/X})^{\sim} \;=\; M^{\sim}_{/X} \;. \]
I would expect definitely this to be the case for each of the categories
(Gr-stacks), (n-Gr-stacks) (although this is not really a model category
in the sense of Quillen,\hand{as} it gives rise only to n-truncated homotopy
types ...),  (stacks)  and (n-stacks),  corresponding to an arbitrary \margin{$\lceil$}
choice of a coherator,  defining the notion of a (Gr-stack),  or of a \margin{$[$}
stack  (for the latter and relations between the two,  see below). \margin{$\lfloor$}

\margin{$|$}\unsure{\over{1}{2}}1.  I would like to digress a little more,  to emphasize still about
\margin{$\downarrow$}the vast variety of algebraic structures giving rise to model categories
for the usual homotopy category,  or at any rate suitable for expressing
more or less arbitrary homotopy types.  Apart from stacks,  where every
thing is still heuristics for the time being,  we have noticed so far
three examples of such structures,  namely  ss complexes,  cubical com-%
plexes,  and categories.  There are a  few familiar variants of the two
former,  such as the  ``simplicial complexes'' (in contrast to semi-simpli
cial ones) ,  namely \inserted[typed]{presheaves on} \xout{functors from} the category of non-empty finite sets\hand{,}
which can be interpreted as ss complexes enriched with symmetry operati-%
ons on each component  -  and there is the corresponding variant in the
cubical case.  It would be rather surprising that these were not just
as good model categories,  as the more habitual ones (\over{*}{16}) - all the more
as the singular complex (simplicial or cubical) is naturally endowed
with this extra structure,  which one generally chooses to forget. More
interesting variants are the  n-multicomplexes (simplicial or cubical,
with or without symmetries),  defined \hand{by} \xout{presheaves} contravariant func-%
tors to sets  in n arguments rather than in just one,  where $n\hand{\geqslant}\hand{1}$. The-%
se complexes are familiar mainly,  it seems,  because of their connecti\hand{on}
with product spaces and the Künneth-Eilenberg-Zilber type relations.
It is generally understood that to such a\xout{n} multicomplex is associated
